\documentclass[11pt]{article}

\usepackage[T1]{fontenc}
\usepackage[margin=1in]{geometry}
\usepackage{amsmath,amssymb,amsthm}
\usepackage{braket}
\usepackage{array}
\usepackage{hyperref}

\newtheorem{theorem}{Theorem}
\newtheorem{lemma}{Lemma}
\newtheorem{proposition}{Proposition}
\theoremstyle{definition}
\newtheorem{definition}{Definition}
\newtheorem{remark}{Remark}

\newcommand{\F}{\mathbb{F}}
\newcommand{\cat}{\mathrm{cat}}

\title{Bounding $\chi(\cat_6\otimes\cat_6)$}
\author{}
\date{}

\begin{document}

\maketitle

\section{Problem statement}

\begin{definition}[Stabilizer state]
An $n$-qubit pure state $\ket{\varphi}$ is a \emph{stabilizer state} if it can be written, in some computational basis, as
\[
\ket{\varphi}=\sum_{x\in A}i^{\ell(x)}(-1)^{q(x)}\ket{x},
\]
where $A\subseteq\F_2^n$ is an affine subspace, $q:A\to\{0,1\}$ is a quadratic function, and $\ell:A\to\{0,1,2,3\}$ is a linear function.
\end{definition}

\begin{definition}[Stabilizer rank]
For a state $\ket{\psi}$, the stabilizer rank $\chi(\psi)$ is the minimum $k$ such that
\[
\ket{\psi}=\sum_{i=1}^k c_i\ket{\varphi_i}
\]
for some stabilizer states $\ket{\varphi_i}$ and coefficients $c_i\in\mathbb{C}$.
\end{definition}

\begin{definition}[The $T$-state]
Define the unnormalized states
\[
\ket{T}:=\ket{0}+e^{i\pi/4}\ket{1},
\qquad
\ket{T^\perp}:=\ket{0}-e^{i\pi/4}\ket{1}.
\]
\end{definition}

\begin{definition}[Cat states]
For a single-qubit state $\ket{\psi}$ with orthogonal complement $\ket{\psi^\perp}$, with phase fixed by $\braket{0|\psi^\perp}$ real and positive, define
\[
\ket{\cat_n(\psi)}:=\frac{1}{\sqrt{2}}
\left(\ket{\psi}^{\otimes n}+\ket{\psi^\perp}^{\otimes n}\right).
\]
Throughout, ``$\cat_6$'' means specifically $\cat_6(T)$:
\[
\ket{\cat_6}:=\frac{1}{\sqrt{2}}
\left(\ket{T}^{\otimes6}+\ket{T^\perp}^{\otimes6}\right).
\]
\end{definition}

It is known that $\chi(\cat_6)\leq3$ (Qassim, Waterloo PhD thesis 2020, \S4.2), so by submultiplicativity
\[
\chi(\cat_6\otimes\cat_6)\leq9.
\]
The problem is to determine $\chi(\cat_6\otimes\cat_6)$, or the best bound that can be proved. In particular, does
\[
\chi(\cat_6\otimes\cat_6)\leq8
\]
hold?

\section{Conventions}

I interpret ``computational basis'' as the fixed standard computational basis. Clifford changes of basis may be used explicitly because stabilizer rank is Clifford-invariant; an arbitrary orthonormal change of basis is not allowed.

A stabilizer state is normally a normalized ray. Throughout the proof I use nonzero, unnormalized representatives of such rays, which does not change stabilizer rank. In the affine-quadratic normal form, choose affine coordinates $t\in\F_2^d$ on the support $A$. The phase is read as
\[
i^{L(t)}(-1)^{Q(t)},
\]
where $L$ is a linear polynomial modulo $4$ and $Q$ is a quadratic polynomial modulo $2$ (constant and linear terms may be absorbed between the two factors). This is the standard meaning of the phase convention in the question.

\section{Best proved bound}

Let $C=\ket{\cat_6}$. The rigorous conclusion is
\[
\boxed{3\leq\chi(C\otimes C)\leq9.}
\]
The upper bound is given by an exact nine-term decomposition below. The lower bound uses the published theorem $\chi(\cat_6)=3$ and a stabilizer contraction. I do not obtain an eight-term decomposition, nor a lower bound excluding rank eight. Therefore this is a progress report, not a complete solution. The question whether
\[
\chi(C\otimes C)\leq8
\]
remains unresolved here; it is Qassim's Conjecture (4.30) in thesis \S4.4.

\section{An exact three-term decomposition of $C$}

Put
\[
\omega=e^{i\pi/4},\qquad
V=\{x\in\F_2^6:|x|\equiv0\pmod2\},\qquad
S(x)=\sum_{1\leq j<k\leq6}x_jx_k\pmod2,
\]
and define the following unnormalized states:
\[
\begin{aligned}
\ket{g}&=\ket{0^6}-i\ket{1^6},\\
\ket{e}&=\sum_{x\in V}\ket{x},\\
\ket{k}&=\sum_{x\in V}(-1)^{S(x)}\ket{x}.
\end{aligned}
\]
All three are stabilizer states in the affine-quadratic sense. For $e$ and $k$, the support is the even-parity linear code and the phases are respectively $1$ and $(-1)^S$. For $g$, parametrize the repetition code by $x=t(1,1,1,1,1,1)$; the phase is
\[
i^t(-1)^t=(-i)^t,
\]
which gives exactly the two displayed amplitudes.

Since $T^\perp=ZT$,
\[
C=\frac1{\sqrt2}\left(T^{\otimes6}+(T^\perp)^{\otimes6}\right)
=\sqrt2\sum_{x\in V}\omega^{|x|}\ket{x}.
\]
I claim that
\begin{equation}
\boxed{
C=2\sqrt2\,\ket{g}
+\frac{-1+i}{\sqrt2}\ket{e}
+\frac{-1-i}{\sqrt2}\ket{k}.
}
\tag{1}
\end{equation}
For an even-weight string $x$,
\[
(-1)^{S(x)}=(-1)^{\binom{|x|}{2}},
\]
whose signs at weights $0,2,4,6$ are $+1,-1,+1,-1$. Set
\[
a=2\sqrt2,\qquad b=\frac{-1+i}{\sqrt2},\qquad
c=\frac{-1-i}{\sqrt2}.
\]
The amplitudes of $ag+be+ck$, grouped by Hamming weight, are
\[
\begin{array}{c|c|c}
|x|&\text{amplitude}&\sqrt2\,\omega^{|x|}\\ \hline
0&a+b+c=\sqrt2&\sqrt2\\
2&b-c=i\sqrt2&i\sqrt2\\
4&b+c=-\sqrt2&-\sqrt2\\
6&-ia+b-c=-i\sqrt2&-i\sqrt2.
\end{array}
\]
Both sides vanish on odd-weight strings, so this proves (1) on every computational-basis vector.

\section{The nine-term upper bound}

Let
\[
d_g=2\sqrt2,\qquad d_e=\frac{-1+i}{\sqrt2},\qquad
d_k=\frac{-1-i}{\sqrt2}.
\]
Tensoring (1) with itself gives the exact identity
\begin{equation}
\boxed{
C\otimes C
=\sum_{u,v\in\{g,e,k\}}d_ud_v\,\ket{u}\otimes\ket{v}.
}
\tag{2}
\end{equation}
The tensor product of two stabilizer states is a stabilizer state: its affine support is the Cartesian product of the two supports, while its linear and quadratic phases are the sums of the phases of the factors. Equation (2) is therefore a nine-term stabilizer decomposition, and
\[
\chi(C\otimes C)\leq9.
\]

For clarity, the nine states in this particular expansion are linearly independent. Indeed, a relation among $g,e,k$ evaluated on strings of weights $2$ and $4$ first forces the coefficients of $e,k$ to vanish, and then the coefficient of $g$ vanishes. Tensor products of two independent triples are independent. This only shows that one cannot delete a term from (2); it is not an obstruction to an eight-term decomposition using different, possibly entangled, stabilizer states.

There is a modest general obstruction to improving (2) while keeping every summand product across the displayed $6|6$ cut. Suppose
\[
C\otimes C=\sum_{j=1}^k c_j\,\ket{\alpha_j}\otimes\ket{\beta_j},
\]
where every $\alpha_j$ and $\beta_j$ is a six-qubit stabilizer state. Group proportional left factors, absorbing proportionality constants into the coefficients, and denote the distinct resulting left factors by $\alpha_{(1)},\ldots,\alpha_{(m)}$. If these vectors are linearly independent, the decomposition has the form
\[
C\otimes C=\sum_{p=1}^m\ket{\alpha_{(p)}}\otimes\ket{u_p}.
\]
Contracting the right block against a linear functional nonzero on $C$ first shows that $C$ belongs to the span of the $\alpha_{(p)}$. Uniqueness of coordinates in this independent left family then gives scalars $\lambda_p$ such that
\[
C=\sum_{p=1}^m\lambda_p\ket{\alpha_{(p)}},
\qquad \ket{u_p}=\lambda_p C.
\]
At least three $\lambda_p$ are nonzero because $\chi(C)=3$. For every such $p$, the terms grouped into $u_p$ express a nonzero multiple of $C$ using the stabilizer states $\beta_j$, and hence that group contains at least three terms. Consequently $k\geq9$. Thus an eight-term decomposition consisting only of product stabilizer states would have to use linearly dependent sets of distinct left factors and, by the same argument, distinct right factors.

The necessary dependence can also be quantified without assuming the distinct factors independent. Regard the decomposition as the amplitude-matrix identity
\[
ADB^{\mathsf T}=CC^{\mathsf T},
\]
where the columns of $A$ and $B$ are the amplitude vectors of $\alpha_j$ and $\beta_j$, and $D=\operatorname{diag}(c_1,\ldots,c_k)$ after zero terms are removed. If $r=\operatorname{rank}A$ and $r'=\operatorname{rank}B$, Sylvester's rank inequality gives
\[
1=\operatorname{rank}(ADB^{\mathsf T})\geq r+r'-k.
\]
Moreover $C$ belongs to both column spans: contract the other tensor factor against any linear functional nonzero on $C$. Since each column is a stabilizer state and $\chi(C)=3$, it follows that $r,r'\geq3$. Any eight-term product-cut decomposition must therefore satisfy
\[
3\leq r,r',\qquad r+r'\leq9.
\]
This argument does not apply to stabilizer states entangled across the $6|6$ cut, which are allowed in the definition of $\chi(C\otimes C)$.

\section{A smaller exact target for an eight-term construction}

Put
\[
b=\frac{-1+i}{\sqrt2},\qquad c=\frac{-1-i}{\sqrt2},\qquad
\ket{D}=b\ket{e}+c\ket{k}.
\]
Thus $C=2\sqrt2\,g+D$, and hence
\begin{equation}
C\otimes C
=8\,g\otimes g+2\sqrt2\,(g\otimes D+D\otimes g)+D\otimes D.
\tag{3}
\end{equation}
Expanding each occurrence of $D$ in the two mixed terms costs two stabilizer states. Consequently
\begin{equation}
\boxed{\chi(C\otimes C)\leq5+\chi(D\otimes D).}
\tag{4}
\end{equation}
In particular, an exact three-stabilizer decomposition of $D\otimes D$ would prove the desired bound $\chi(C\otimes C)\leq8$.

This reduced target has a simple five-qubit description, but the compression must be implemented as a six-bit reversible map. Let $U$ be the computational-basis permutation induced by
\[
(x_1,\ldots,x_6)\longmapsto(p,y_1,\ldots,y_5)
=(x_1+\cdots+x_6,x_2,\ldots,x_6).
\]
Its inverse is
\[
x=(p+y_1+\cdots+y_5,y_1,\ldots,y_5),
\]
so $U$ is an invertible linear transformation implemented by CNOT gates. On the even-parity support $V$, the first output bit is $p=0$. If
\begin{equation}
q(y)=\sum_{1\leq j<k\leq5}y_jy_k+\sum_{j=1}^5y_j\pmod2,
\tag{5}
\end{equation}
then $S(x)=q(y)$ when $p=0$. Since the amplitudes of $D$ are $-\sqrt2$ for $S=0$ and $i\sqrt2$ for $S=1$, the exact transformed state is
\begin{equation}
\ket{d}=\sum_{y\in\F_2^5}i^{-q(y)}\ket{y},
\qquad U\ket{D}=-\sqrt2\,\ket{0}\otimes\ket{d}.
\tag{6}
\end{equation}
Here $q(y)$ in the exponent is represented by the integer $0$ or $1$. Clifford invariance, adjoining a stabilizer ancilla, and contraction with $\bra{0}$ give the two directions needed for
\[
\chi(D)=\chi(d),\qquad \chi(D\otimes D)=\chi(d\otimes d).
\]
Accordingly, the missing certificate in this route is a three-term affine-quadratic decomposition of the ten-qubit phase function
\begin{equation}
(y,z)\longmapsto i^{-q(y)-q(z)}.
\tag{7}
\end{equation}

The state $D$ really has rank two, rather than one. It already has the two-term decomposition $be+ck$. On the other hand, a full-support stabilizer phase has an exponent modulo four of the form $L+2Q$, so its mixed second difference is even. For the exponent $h=-q\pmod4$ of (6), evaluated at $0,e_1,e_2,e_1+e_2$, the values are $0,3,3,3$, and therefore
\[
\Delta_{e_1}\Delta_{e_2}h(0)=3-3-3+0=1\pmod4.
\]
Thus (6) is not a stabilizer state and $\chi(D)=2$.

There are exactly two stabilizer rays in the plane spanned by the two full-support stabilizers $1$ and $(-1)^q$. Indeed, up to overall scale, any vector in that plane has one amplitude $a+b$ on $q=0$ and another $a-b$ on $q=1$. If it is a full-support stabilizer state, their ratio must lie in $\{1,-1,i,-i\}$. Ratios $1$ and $-1$ give the original two rays. Ratios $i$ and $-i$ give phases proportional to $i^q$ and $i^{-q}$, and the same mixed-second-difference calculation excludes both. If one amplitude vanishes, the support is a level set of $q$. Translation by $r=1^5$ interchanges the two level sets because $q(y+r)=q(y)+1$; hence each has size $2^4=16$. The zero level is not affine: it contains $0$, $u=e_1+e_2+e_3$, and $v=e_1+e_2+e_4$, but not $u+v=e_3+e_4$. The one level cannot be affine either, because its complement would then be the opposite affine hyperplane. Hence no further stabilizer ray occurs.

This plane rigidity has a tensor-square analogue. Write
\[
\ket{u}=\sum_y\ket{y},\qquad
\ket{v}=\sum_y(-1)^{q(y)}\ket{y},
\]
so that $u,v$ are orthogonal stabilizer states and
\[
\ket{d}=\frac{1-i}{2}(\ket{u}+i\ket{v}).
\]
The only stabilizer rays in
\[
\operatorname{span}\{u\otimes u,u\otimes v,v\otimes u,v\otimes v\}
\]
are the four displayed product rays. To prove this, observe that every vector in this four-dimensional space has a constant amplitude $m_{ab}$ on each cell
\[
\{(y,z):q(y)=a,\ q(z)=b\},\qquad (a,b)\in\F_2^2,
\]
and every cell has size $16^2=256$. If its support is affine, the number of nonzero cells is therefore $1$, $2$, or $4$. A one-cell support is not affine, because its projection onto either block is a nonaffine level set of $q$. A two-cell support is either a row, a column, or one of the two diagonals. Rows and columns again have a nonaffine projection. A diagonal is a level set of the balanced quadratic $q(y)+q(z)$; if it were an affine hyperplane, that Boolean quadratic would be affine, contrary to its nonzero polar form.

It remains to consider full support. Restricting a putative stabilizer state to a fixed computational-basis value of $z$ gives a full-support stabilizer in $\operatorname{span}\{u,v\}$. The preceding plane result forces the ratio $m_{1b}/m_{0b}$ to be $\pm1$, and restricting in the other direction gives the analogous conclusion for rows. After an overall scaling, the table $(m_{ab})$ is therefore a sign table. Every two-by-two sign table is either a product of a sign depending only on $a$ and one depending only on $b$, or differs from such a product by $(-1)^{ab}$. The product tables give exactly the four product rays above. The remaining possibility is not stabilizer: its exponent contains $2q(y)q(z)$ modulo four, whose fourth additive difference in the directions
\[
(e_1,0),\ (e_2,0),\ (0,e_1),\ (0,e_2)
\]
at the origin is $2\pmod4$. Every stabilizer exponent $L+2Q$ is quadratic modulo four and has vanishing fourth additive differences. This proves the claim. Since all four coordinates of $d\otimes d$ in this product basis are nonzero, no three stabilizer rays lying in this four-dimensional plane can span it. Thus any three-term decomposition sought in (7) must use at least one stabilizer state outside the natural tensor-square plane.

There is at least one exact entangled stabilizer slice in (7). The polar form of $q$ is
\[
B(y,z)=\left(\sum_jy_j\right)\left(\sum_jz_j\right)+y\mathbin{\cdot} z,
\]
whose radical is generated by $r=1^5$, and $q(y+r)=q(y)+1$. Hence (7) is the constant $-i$ on the affine graph $z=y+r$. The uniform state on that graph is a stabilizer state entangled across the $5|5$ cut. This observation has not been extended here to a three-term decomposition of the full phase function (7); the ordinary four products of the two summands of $D$ remain the only certified decomposition. Thus (4) currently reproduces nine, not eight.

The graph cannot be enlarged in the most obvious maximal-isotropic way. Let
\[
L=\{(y,z):z=y+sr,\ y\in\F_2^5,\ s\in\F_2\}.
\]
Its six-dimensional direction space is totally isotropic for $B\oplus B$, but that parity condition is only necessary, not sufficient, for the restriction of (7) to be a stabilizer phase. Indeed, on $L$ the target phase is $(-1)^{q(y)}$ for $s=0$ and $-i$ for $s=1$. One exponent modulo four is
\[
H(y,s)=2q(y)+3s+2s q(y)\pmod4.
\]
Its third additive difference in the directions $e_1,e_2$ of the $y$ block and the $s$ direction is $2\pmod4$, whereas every stabilizer exponent has vanishing third additive differences. Hence this 64-point restriction is not a stabilizer state; only the 32-point graph slice above is certified.

There is, however, a finite support-dimension reduction for deciding whether the missing three-term decomposition exists. Let
\[
\ket{\psi}=\ket{d}\otimes\ket{d}
\]
and suppose $\psi=\sum_{j=1}^3c_j\varphi_j$, where $A_j\subseteq\F_2^{10}$ is the affine support of $\varphi_j$. Because $\psi$ has full support, the $A_j$ cover $\F_2^{10}$; in particular, at least one has dimension at least nine, since $3\cdot2^8<2^{10}$.

More strongly, either all three $A_j$ have full dimension, or some affine hyperplane-coset restriction of $\psi$ has stabilizer rank at most two. To see this, if $A_j$ is proper, place it inside an affine hyperplane $H$ and let $K=\F_2^{10}\setminus H$, which is the opposite coset. Projection onto $K$ is a Pauli-$Z$ eigenspace projection, so it takes every stabilizer state to zero or a stabilizer state. It annihilates $\varphi_j$, leaving at most two summands for $\psi|_K$.

In fact, every one of these hyperplane-coset restrictions has rank at least two. The exponent of (7), reduced modulo two, is $q(y)+q(z)$, with polar form $B\oplus B$. The radical of this alternating form has dimension two, and its nondegenerate quotient has dimension eight, so a totally isotropic subspace has dimension at most $2+8/2=6$. If the restriction to an affine hyperplane were a stabilizer state, its nine-dimensional direction space would have to be totally isotropic: otherwise some mixed second difference of its exponent would be odd. This is impossible. Hence the restrictions arising above would have rank exactly two, not one.

Consequently, proving $\chi(d\otimes d)=4$ can be reduced to two exact tasks: exclude rank two for each of the $2(2^{10}-1)=2046$ hyperplane-coset restrictions, and exclude representations by at most three full-support stabilizer phases. Repeated rays are allowed in this statement (and may of course be combined); thus the second task includes possible ranks one and two, not only decompositions having three distinct rays. This is a finite reduction, not its completion; no classification or exhaustive certificate carrying out those two tasks is supplied here. It also constrains the five-dimensional graph slice: if it were one summand and the other two supports were both proper, those two supports would have to be the two complementary affine hyperplanes (otherwise their union has at most $768<992$ points), although a full-support remaining summand is also possible.

\section{The lower bound}

Qassim, Pashayan, and Gosset prove in the appendix of \emph{Improved Upper Bounds on the Stabilizer Rank of Magic States} that
\begin{equation}
\chi(\cat_6)=3.
\tag{8}
\end{equation}
The paper uses $\ket{T}_{\mathrm{norm}}=2^{-1/2}(\ket{0}+\omega\ket{1})$, whereas the problem uses the unnormalized vector. Its cat state is consequently $C/8$. Multiplication by a nonzero overall scalar does not affect stabilizer rank, so (8) applies with the convention here.

Their proof first shows $\chi(\cat_5)>2$ by a canonical-form classification of pairs of stabilizer states and a finite Pauli-spectrum comparison, and then uses the computational-basis contraction $\braket{0|\cat_6}\propto\cat_5$. I use (8) as a cited theorem rather than claiming a self-contained reproduction of that classification and computer-assisted finite check.

Now contract the second six-qubit block of $C\otimes C$ with the stabilizer bra $\bra{0^6}$. Since
\[
\braket{0^6|C}=\sqrt2,
\]
we obtain
\begin{equation}
(I\otimes\bra{0^6})(C\otimes C)=\sqrt2\,C.
\tag{9}
\end{equation}
Contracting a stabilizer state with a computational-basis stabilizer bra gives either zero or another stabilizer state. In affine coordinates, one intersects the affine support with the equations fixing the contracted coordinates and restricts the same affine-quadratic phase to that intersection. Hence an $r$-term decomposition of $C\otimes C$ would induce an at-most-$r$-term decomposition of $C$. Equations (8) and (9) imply
\[
\chi(C\otimes C)\geq3.
\]

\section{A further exact contraction}

There is a useful further reduction of the lower-bound problem. Let
\[
\ket{B}=\ket{00}+i\ket{11},
\]
which is a two-qubit stabilizer state, and contract its bra against one qubit from each copy of $C$. Writing $U=T^\perp$, direct calculation gives
\[
\braket{B|T,T}=\braket{B|U,U}=2,
\qquad
\braket{B|T,U}=\braket{B|U,T}=0.
\]
Therefore, up to the harmless ordering of the ten uncontracted qubits,
\begin{equation}
\bra{B}\,C\otimes C
=T^{\otimes10}+U^{\otimes10}
=\sqrt2\,\ket{\cat_{10}}.
\tag{10}
\end{equation}
For completeness, choose a two-qubit Clifford $W$ such that $W\ket{00}\propto\ket{B}$. Then $\bra{B}\propto\bra{00}W^\dagger$. Thus contraction with $\bra{B}$ is a Clifford operation on the contracted qubits followed by computational-basis restriction, and it maps each stabilizer summand to zero or a stabilizer state. Therefore
\begin{equation}
\chi(C\otimes C)\geq\chi(\cat_{10}).
\tag{11}
\end{equation}
This is precisely the $\cat_2$ contraction used by Qassim, Pashayan, and Gosset to construct larger cats. At present it does not strengthen the numerical lower bound: computational-basis contractions show $\chi(\cat_{10})\geq\chi(\cat_6)=3$, while no proof of $\chi(\cat_{10})\geq4$ is supplied here. Equation (11) nevertheless identifies a concrete smaller exact-rank problem whose solution would immediately improve the lower bound.

There is also an exact relation in the other direction. Let
\[
\Pi=\frac12\left(I+Z^{\otimes6}\otimes I^{\otimes6}\right).
\]
The state $\cat_{12}$ is supported on strings of even total parity. Projecting its first six qubits onto even parity therefore forces each block to have even parity. Comparing computational-basis amplitudes gives
\begin{equation}
\boxed{C\otimes C=\sqrt2\,\Pi\ket{\cat_{12}}.}
\tag{12}
\end{equation}
Pauli projection maps a stabilizer state to zero or a stabilizer state, so together with (11) this gives the rigorous sandwich
\begin{equation}
\chi(\cat_{10})\leq\chi(C\otimes C)
\leq\chi(\cat_{12}).
\tag{13}
\end{equation}
No bound on $\cat_{12}$ strong enough to improve nine is established here.

\section{Status and remaining open issues}

\subsection*{What has been rigorously established}

The proved numerical result is exactly
\[
\boxed{3\leq\chi(C\otimes C)\leq9.}
\]
The upper bound is certified by the explicit three-term decomposition (1) of $C$ and its nine-term tensor square (2). The lower bound follows from the published exact theorem $\chi(\cat_6)=3$ and the contraction (9). The additional exact relations
\[
\chi(\cat_{10})\leq\chi(C\otimes C)\leq\chi(\cat_{12})
\]
are also established, but they do not improve the numerical interval here.

For product stabilizer summands across the natural $6|6$ cut, the argument above proves that an eight-term decomposition with independent distinct left factors is impossible, and likewise with left and right exchanged. More generally, any eight-term product-cut decomposition must have local ranks $r,r'\geq3$ satisfying $r+r'\leq9$, so both local collections must have the necessary dependencies described above. These restrictions do not apply to genuinely entangled stabilizer summands.

The concrete upper-bound route through (4) is also rigorously reduced. The state $D$ has rank exactly two and is Clifford-equivalent, up to a stabilizer ancilla and a nonzero scalar, to the five-qubit state $d$ in (6). A three-term stabilizer decomposition of $d\otimes d$ would yield the desired eight-term decomposition. Such a decomposition cannot remain inside the natural four-dimensional tensor-square plane, whose only stabilizer rays are the four product rays. The affine graph $z=y+1^5$ gives one exact entangled stabilizer slice, while its proposed 64-point maximal-isotropic enlargement is rigorously excluded by the third-difference calculation. Finally, exclusion of a three-term decomposition of $d\otimes d$ is reduced to excluding rank two for all $2046$ hyperplane-coset restrictions and excluding representations by at most three full-support stabilizer phases. No hyperplane-coset restriction has rank one.

\subsection*{What remains open}

\begin{enumerate}
\item \emph{Eight-term upper bound.} The gap is whether the upper bound after (2) can be lowered to eight. The exact product construction (2) was analyzed, including the independence of its nine summands. The product-cut argument also forces strong dependencies in any eight-term decomposition of product stabilizer states, but it does not constrain genuinely entangled stabilizer summands. Equation (4) reduces one concrete route to finding a three-term decomposition of the ten-qubit phase function (7). The tensor-square-plane classification proves that such a decomposition cannot stay inside the natural span of the four product summands, and the tempting 64-point enlargement of the known graph slice was excluded by a third-difference obstruction. The support-dimension argument further reduces exclusion of this route to hyperplane restrictions and the full-support case, but those finite checks were not completed. Either a three-term certificate for (7), necessarily using stabilizers outside the natural tensor-square plane, or an affine-quadratic description of eight different stabilizer summands with an exact amplitude verification, would close the upper-bound gap.

\item \emph{Lower bound above three.} The gap in (9)--(13) is that the contractions certify only rank three. The $\cat_2$ contraction (10) reduces to $\cat_{10}$, but the cited small-cat classification does not prove $\chi(\cat_{10})\geq4$; the projection from $\cat_{12}$ also supplies no stronger number. Excluding all three-stabilizer spans containing $\cat_{10}$, or finding another contraction with a known rank-four output, would close this gap; it would improve the interval but would not settle the eight-term conjecture.

\item \emph{Exact value.} The final gap is that ranks $3,4,\ldots,9$ remain compatible with the upper decomposition and lower contractions. Those are the two proof mechanisms used here, and neither supplies matching bounds. An exact $k$-term witness together with an exclusion of all ranks below $k$ would determine the value.
\end{enumerate}

\subsection*{Approaches tried and why they stalled}

The exact product construction was analyzed first. Its nine displayed summands are independent, so no term can simply be deleted. The product-cut arguments impose dependencies on any hypothetical eight-term product decomposition, but they do not constrain stabilizer states entangled across the cut, which remain allowed.

The reduced target $D\otimes D$, or equivalently $d\otimes d$, was then analyzed. Its ordinary product expansion has four terms and therefore only recovers the nine-term bound. The natural tensor-square plane contains no other stabilizer rays, so a three-term decomposition cannot be found wholly within that plane. The known entangled graph slice did not extend to the tempting 64-point support because the restricted exponent has a nonzero third additive difference. The remaining support classification was reduced to finite exact tasks, but those tasks were not completed.

The lower-bound approach by contraction reached $\cat_{10}$, but the available pair-classification result only supplies the threshold three for the cited small-cat problem. Proving rank at least four would require control of spans of three stabilizer states, to which the pair normal form does not directly generalize. The projection from $\cat_{12}$ similarly supplies no stronger numerical upper bound.

Computational investigations were also tried as search guidance, but none yielded an exact certificate:
\begin{itemize}
\item The five-qubit Pauli-spectrum enumeration gave
\[
\mathrm{PS}(\cat_5):
(-1/2)^{20},(-1/4)^{40},0^{782},
(1/4)^{120},(1/2)^{60},1^2.
\]
Orthogonal canonical rank-two candidates had $710$ zero Pauli expectations. The nonorthogonal candidates $\gamma=\pm i$ had $722$ zeros. The candidate $\gamma=-1/2$ had $782$ zeros but thirty expectations of absolute value $3/4$. This finite comparison is part of the cited small-cat line of attack, but it did not provide control of spans of three stabilizer states or prove $\chi(\cat_{10})\geq4$.

\item The Clifford $A=e^{-i\pi/4}SX$ has eigenvectors $T,T^\perp$ with eigenvalues $+1,-1$. Applying $A$ on any even subset of the six qubits fixes $C$. Applying these symmetries to $g,e,k$ generated $48$ distinct stabilizer rays spanning a $32$-dimensional subspace. Randomized searches over products of these rays found no exact eight-term witness. This is evidence only about that restricted pool.

\item The parity support was Clifford-compressed using the six-bit invertible map
\[
x\longmapsto(x_1+\cdots+x_6,x_2,\ldots,x_6).
\]
Thus one cat state is Clifford-equivalent to a five-qubit state tensored with a stabilizer ancilla, and its square reduces in the same way to a ten-qubit target. A simulated-annealing search at rank eight did not yield a certificate; this is not evidence against the conjecture.

\item A warm-started simulated-annealing search for a rank-three decomposition of the smaller target $D^{\otimes2}$ did not improve on the residual of a leave-one-out four-product seed. This is search evidence only, not a lower bound on $\chi(D^{\otimes2})$.
\end{itemize}

Other possible shortcuts do not close the problem. The generic multiplicativity result of Lovitz--Steffan does not decide this special algebraic rank-two state. Stabilizer fidelity or extent has not been used as a rank lower bound because the needed general inequality between extent and rank is not supplied by the usual Cauchy--Schwarz argument for an ill-conditioned minimal stabilizer spanning set. The 2026 Labib--Russo four-copy qubit identity concerns a different Clifford orbit and therefore does not provide the required decomposition for the phase state used here.

\subsection*{Promising next steps}

For the upper bound, a successful next step would be an exact affine-quadratic search over genuinely entangled ten- or twelve-qubit stabilizer states. Concretely, it could produce a three-term certificate for the phase function (7), necessarily using at least one stabilizer outside the natural tensor-square plane, or directly produce eight stabilizer summands for $C\otimes C$ together with exact tableau data, exact algebraic coefficients, and a computational-basis amplitude verification. For exclusion along the reduced route, it would suffice to complete the exact classification of the $2046$ rank-two hyperplane-coset restrictions and the at-most-three full-support phase representations.

For the lower bound, a promising next step would be a classification of three-stabilizer spans sufficient to decide whether $\chi(\cat_{10})\geq4$, or a different stabilizer contraction whose output has a known rank-four lower bound. Such a result would improve the interval, though it would not by itself decide the eight-term conjecture. Determining the exact value ultimately requires both an exact $k$-term witness and an exclusion of all smaller ranks.

\begin{thebibliography}{9}

\bibitem{QassimThesis}
Hammam Qassim,
\emph{Classical Simulations of Quantum Systems Using Stabilizer Decompositions},
PhD thesis, University of Waterloo, 2020; deposited in UWSpace in 2021.
\url{https://dspacemainprd01.lib.uwaterloo.ca/server/api/core/bitstreams/aff1f984-003a-4801-bd9c-608f0bfb8743/content}.
The official repository record, dated February 19, 2021, is
\url{https://uwspace.uwaterloo.ca/items/147b4d1e-1a97-49a8-a3e0-f5954c570cdb}.
The conjecture $\chi(\cat_6^{\otimes2})\leq8$ is equation (4.30) in \S4.4, ``Discussion, conjectures, and open problems''; the conjectural exponent $3/8$ is equation (4.31).

\bibitem{QPG2021}
Hammam Qassim, Hakop Pashayan, and David Gosset,
\emph{Improved Upper Bounds on the Stabilizer Rank of Magic States},
\emph{Quantum} \textbf{5}, 606 (2021).
\url{https://quantum-journal.org/papers/q-2021-12-20-606/}.
The three-term $\cat_6$ construction appears in the main text, and the appendix proves the exact value $\chi(\cat_6)=3$.

\bibitem{BSS2016}
Sergey Bravyi, Graeme Smith, and John A. Smolin,
\emph{Trading Classical and Quantum Computational Resources},
\emph{Physical Review X} \textbf{6}, 021043 (2016).
\url{https://doi.org/10.1103/PhysRevX.6.021043}.
Background on exact stabilizer rank and magic-state simulation.

\bibitem{BBC2019}
Sergey Bravyi, David Browne, Padraic Calpin, Earl Campbell, David Gosset, and Mark Howard,
\emph{Simulation of Quantum Circuits by Low-Rank Stabilizer Decompositions},
\emph{Quantum} \textbf{3}, 181 (2019).
\url{https://quantum-journal.org/papers/q-2019-09-02-181/}.
Background on stabilizer decompositions, stabilizer fidelity, and stabilizer extent.

\bibitem{LovitzSteffan2022}
Benjamin Lovitz and Vincent Steffan,
\emph{New Techniques for Bounding Stabilizer Rank},
\emph{Quantum} \textbf{6}, 692 (2022).
\url{https://doi.org/10.22331/q-2022-04-20-692}.
Besides general lower-bound techniques, this paper constructs rank-two states whose tensor squares have rank four. Its generic multiplicativity result does not decide the special algebraic state $D$ isolated here.

\end{thebibliography}

\end{document}
